\documentclass[10pt, spanish]{beamer}

% --- PAQUETES ---
\usetheme{metropolis}           
\usepackage[utf8]{inputenc}
\usepackage[spanish]{babel}
\usepackage{listings}
\usepackage{xcolor}
\usepackage{graphicx}
\usepackage{amsmath, amssymb}
\usepackage{karnaugh-map} % Para mapas de Karnaugh
\usepackage{tikz}         % Para diagramas y circuitos
\usepackage{circuitikz}   % Para circuitos eléctricos
\usepackage{booktabs}     % Para tablas elegantes
\usepackage{hyperref}
\usepackage[backend=biber]{biblatex} % Configura Biber como motor
\addbibresource{bibliography.bib}    % Nombre exacto de tu archivo .bib

% --- CONFIGURACIÓN DE CÓDIGO (LISTINGS) ---
\definecolor{codegreen}{rgb}{0,0.6,0}
\definecolor{codegray}{rgb}{0.5,0.5,0.5}
\definecolor{backcolour}{rgb}{0.95,0.95,0.92}

\lstset{
    backgroundcolor=\color{backcolour},   
    commentstyle=\color{codegreen},
    keywordstyle=\color{blue},
    numberstyle=\tiny\color{codegray},
    stringstyle=\color{red},
    basicstyle=\ttfamily\tiny,
    breaklines=true,
    language=[LaTeX]TeX,
    frame=single
}

\title{\LaTeX{} para Ingeniería y Academia}
\subtitle{De la Estructura APA/IEEE}
\author{Introducción a la Composición Tipográfica Profesional}
\date{\today}

\begin{document}

\maketitle

\section{Introducción}

\begin{frame}{¿Qué es \LaTeX? y ¿Por qué usarlo?}
    \begin{itemize}
        \item \textbf{Concepto:} No es un procesador de texto (como Word), es un lenguaje de marcado para composición tipográfica (WYSIWYM).
        \item \textbf{Ventajas:}
            \begin{itemize}
                \item Manejo automático de bibliografía e índices.
                \item Calidad matemática insuperable.
                \item Estabilidad en documentos largos (tesis de 300+ páginas).
                \item Separación total de \textbf{Contenido} y \textbf{Formato}.
            \end{itemize}
    \end{itemize}
\end{frame}

\section{Formatos Estándar: APA 7 e IEEE}

\begin{frame}[fragile]{Estructura APA 7ma Edición}
    Se usa el paquete \texttt{apa7}. Sus modos principales son:
    \begin{itemize}
        \item \textbf{man}: (manuscript) Doble espacio, una columna (para revisión).
        \item \textbf{stu}: (student) Formato para trabajos estudiantiles.
        \item \textbf{jou}: (journal) Formato de revista profesional.
    \end{itemize}
\begin{lstlisting}
\documentclass[stu, spanish]{apa7}
\title{Analisis de Circuitos}
\shorttitle{APA7 en Beamer}
\authorsname{Juan Perez}
\affiliation{Universidad X}
\begin{document}
    \maketitle
    \section{Introduccion} ...
\end{document}
\end{lstlisting}
\end{frame}

\begin{frame}[fragile]{Estructura IEEE}
    Ideal para artículos de ingeniería. Clases comunes:
    \begin{itemize}
        \item \textbf{conference}: Para ponencias en congresos.
        \item \textbf{journal}: Para publicaciones en revistas IEEE.
        \item \textbf{comsoc}: Para comunicaciones.
    \end{itemize}
\begin{lstlisting}
\documentclass[conference]{IEEEtran}
\title{Implementacion de Logica Digital}
\author{\IEEEauthorblockN{Maria Garcia} \IEEEauthorblockA{Escuela de Ingenieria}}
\begin{document}
    \maketitle
    \begin{abstract} Resumen del articulo... \end{abstract}
\end{document}
\end{lstlisting}
\end{frame}

\section{Elementos Académicos Comunes}

\begin{frame}[fragile]{Ecuaciones y Referencias}
    Diferencias fundamentales en matemáticas:
    \begin{itemize}
        \item \textbf{En línea:} \lstinline|$E=mc^2$| para el texto ($E=mc^2$).
        \item \textbf{Bloque:} \lstinline|$$...$$| o \lstinline|\[...\]| no se numeran.
        \item \textbf{Entorno \texttt{equation}:} Se numera y permite usar \texttt{\textbackslash label}.
    \end{itemize}
\begin{lstlisting}
\begin{equation}
    V = I \cdot R \label{eq:ohm}
\end{equation}
Como se observa en la ec. (\ref{eq:ohm})...
\end{lstlisting}
    \textbf{Salida:}
    \begin{equation}
        V = I \cdot R \label{eq:ohm_example}
    \end{equation}
\end{frame}

\begin{frame}[fragile]{Símbolos y Letras Griegas}
    Para símbolos especiales, visita: \href{https://manualdelatex.com/simbolos#chapter3}{Manual de Símbolos \LaTeX}.
    \begin{itemize}
        \item Letras: \lstinline|$\alpha, \beta, \gamma, \Delta, \Omega$| $\rightarrow$ $\alpha, \beta, \gamma, \Delta, \Omega$.
        \item Operadores: \lstinline|$\int, \sum, \infty, \nabla$| $\rightarrow$ $\int, \sum, \infty, \nabla$.
    \end{itemize}
    Para fórmulas complejas sin escribir código: \href{https://editor-codecogs-com.translate.goog/?_x_tr_sl=en&_x_tr_tl=es&_x_tr_hl=es&_x_tr_pto=tc}{Editor Visual CodeCogs}.
\end{frame}

\begin{frame}[fragile]{Tablas de Cálculo de Errores}
    Recomiendo usar: \href{https://www.tablesgenerator.com/}{Tables Generator}.
\begin{lstlisting}
\begin{table}[h]
\centering
\caption{Comparativa de tensiones.}
\label{tab:errores}
\begin{tabular}{|c|c|c|c|}
\hline
\textbf{Nodo} & \textbf{Medido (V)} & \textbf{Simulado (V)} & \textbf{Error (\%)} \\ \hline
1 & 4.98 & 5.00 & 0.4\% \\ \hline
2 & 2.45 & 2.50 & 2.0\% \\ \hline
\end{tabular}
\end{table}
\end{lstlisting}

También se pueden apoyar de la IA para generar el código de tablas más complejas.
\end{frame}

\begin{frame}[fragile]{Inserción de Imágenes y Figuras}
Para incluir imágenes se usa el paquete \texttt{graphicx} y el entorno \texttt{figure}.
\begin{lstlisting}
\begin{figure}[htbp] % Opciones de posicionamiento
    \caption{Diagrama del sistema propuesto.} % Caption ARRIBA
    \centering
    \includegraphics[width=0.5\textwidth]{mi_imagen.png}
    \label{fig:diagrama} % Etiqueta para referenciar
    \par\vspace{2pt}\small \textit{Nota: Datos obtenidos en laboratorio.} % Nota abajo
\end{figure}
Como se aprecia en la Fig. \ref{fig:diagrama}, los resultados...
\end{lstlisting}

\textbf{Opciones de Posicionamiento [!htbp]:}
\begin{itemize}
    \item \textbf{h} (here): Intenta ponerla justo donde está el código.
    \item \textbf{t} (top): Al principio de la página actual.
    \item \textbf{b} (bottom): Al final de la página actual.
    \item \textbf{p} (page): En una página dedicada solo a flotantes.
    \item \textbf{!} : Fuerza a \LaTeX{} a ignorar algunas reglas estéticas.
\end{itemize}
\end{frame}

\begin{frame}[fragile]{Citas en el Texto: APA 7 vs. IEEE}
La forma de citar varía drásticamente según el estándar:

\begin{table}[]
\tiny
\begin{tabular}{@{}lll@{}}
\toprule
\textbf{Tipo de Cita} & \textbf{Normas APA 7ma Ed.} & \textbf{Normas IEEE} \\ \midrule
\textbf{Parentética} & (García \& Pérez, 2023) & [1] \\
\textbf{Narrativa} & García y Pérez (2023) afirman... & García y Pérez en [1]... \\
\textbf{Múltiple} & (Altamirano, 2020; Zappa, 2018) & [2], [5] \\
\textbf{Páginas} & (García, 2023, p. 45) & [1, p. 45] \\ \bottomrule
\end{tabular}
\end{table}

\textbf{Comandos en \LaTeX:}
\begin{itemize}
    \item \lstinline|\cite{key}|: Cita estándar ([1] o (García, 2023)).
    \item \lstinline|\textcite{key}|: Cita narrativa (García (2023)). Solo en APA/BibLaTeX.
    \item \lstinline|\citeonline{key}|: Variante para IEEE según algunos paquetes.
\end{itemize}
\end{frame}


\begin{frame}[fragile]{Comandos de Citación: \texttt{\textbackslash parentcite} vs \texttt{\textbackslash textcite}}
Para documentos en \textbf{APA 7}, no basta con usar \texttt{\textbackslash cite}. \texttt{biblatex} ofrece comandos específicos para cumplir con la norma:

\begin{itemize}
    \item \textbf{\textbackslash parentcite\{key\}}: Cita tradicional al final de un párrafo.
    \begin{itemize}
        \item \textit{Resultado APA:} (García \& Pérez, 2024)
        \item \textit{Resultado IEEE:} [1]
    \end{itemize}
    \item \textbf{\textbackslash textcite\{key\}}: Cita narrativa donde el autor es el sujeto.
    \begin{itemize}
        \item \textit{Resultado APA:} García y Pérez (2024)
        \item \textit{Resultado IEEE:} García y Pérez [1]
    \end{itemize}
\end{itemize}

\textbf{Ejemplo de uso en el código:}
\begin{lstlisting}
Segun \textcite{manual_latex}, los simbolos son ...
Esto se debe a la estructura del nucelo \parencite{codecogs}.
\end{lstlisting}

\textit{Nota: Usar estos comandos hace que tu documento sea ``estilo-independiente''; si cambias de APA a IEEE, \LaTeX{} ajustará el formato automáticamente.}
\end{frame}


\begin{frame}[fragile]{Configuración de Bibliografía Externa}
Para gestionar referencias, se recomienda el paquete \texttt{biblatex}. En el \textbf{preámbulo} (antes de \texttt{\textbackslash begin\{document\}}) debes agregar:

\begin{lstlisting}
% 1. Cargar el paquete y definir el estilo (apa, ieee, etc.)
\usepackage[style=apa, backend=biber]{biblatex}

% 2. Vincular el archivo .bib (IMPORTANTE: incluir extension)
\addbibresource{referencias.bib}
\end{lstlisting}

\textbf{Puntos clave:}
\begin{itemize}
    \item \textbf{style:} Define si se verá como (García, 2024) o [1].
    \item \textbf{backend=biber:} Es el motor que procesa los datos (más potente que el antiguo BibTeX).
    \item \textbf{Comando de salida:} Usa \texttt{\textbackslash printbibliography} donde quieras que aparezca la lista final.
\end{itemize}
\end{frame}


\begin{frame}[fragile]{Estructura del archivo \texttt{referencias.bib}}
El archivo bibliográfico es una base de datos en texto plano. Cada entrada tiene un tipo (\texttt{@online}, \texttt{@article}, etc.) y una \textbf{llave} única.

\begin{lstlisting}
@online{manual_latex,
    author = {Manual de LaTeX},
    title  = {Simbolos Matematicos},
    year   = {2024},
    url    = {https://manualdelatex.com/simbolos},
    urldate = {2024-05-20}
}

@online{codecogs,
    author = {CodeCogs},
    title  = {Editor Visual de Ecuaciones},
    year   = {2024},
    url    = {https://editor.codecogs.com}
}

@online{tablesgen,
    author = {Tables Generator},
    title  = {Generador de tablas para LaTeX},
    url    = {https://www.tablesgenerator.com}
}
\end{lstlisting}
\textit{Nota: La "llave" (ej. \texttt{codecogs}) es lo que usas para citar con \texttt{\textbackslash cite\{codecogs\}}.}
\end{frame}



\section{Lógica Digital y Gráficos}

\begin{frame}[fragile]{Mapas de K. con Minterminos y Condiciones ``No Importa''}
Podemos optimizar el llenado con \texttt{\textbackslash minterms}, \texttt{\textbackslash maxterms} y usar \textbf{X} (\texttt{\textbackslash autoterms\{X\}}) para simplificar.

\begin{columns}
\begin{column}{0.5\textwidth}
\begin{lstlisting}
\begin{karnaugh-map}[4][4][1][$CD$][$AB$]
	\minterms{4,10,12,14}
	\maxterms{2,3,6,8}
	\autoterms[X]
	\implicant{4}{12}
	\implicant{14}{10}
\end{karnaugh-map}
\end{lstlisting}
\end{column}
\begin{column}{0.5\textwidth}
    \centering
	\begin{karnaugh-map}[4][4][1][$CD$][$AB$]
    		\minterms{4,10,12,14}
	    \maxterms{2,3,6,8}
    		\autoterms[X]
	    \implicant{4}{12} % Grupo de las 4 esquinas
	    \implicant{14}{10}  % Grupo de 2 (celdas 5 y 7)
	\end{karnaugh-map}
\end{column}
\end{columns}

%\small \textbf{Nota:} \texttt{\textbackslash implicant\{0\}\{10\}} en un mapa de 4x4 identifica que los extremos son adyacentes y crea los arcos en las esquinas.
\end{frame}

\begin{frame}[fragile]{El poder de TikZ y Circuitikz}
    Permite dibujar circuitos directamente en el código.
\begin{lstlisting}
\begin{circuitikz}
  \draw (0,0) to[V=$V_s$] (0,2) to[R=$R_1$] (2,2) 
        to[C=$C_1$] (2,0) -- (0,0);
\end{circuitikz}
\end{lstlisting}
    \centering
    \begin{circuitikz}[scale=0.7]
      \draw (0,0) to[V=$V_s$] (0,2) to[R=$10\Omega$] (2,2) 
            to[C=$1\mu F$] (2,0) -- (0,0);
    \end{circuitikz}
\end{frame}

\section{Instalación y Compilación}

\begin{frame}{Compilación Local y Herramientas}
    Para trabajar sin internet necesitas:
    \begin{enumerate}
        \item \textbf{Un Motor (\LaTeX{} Distribution):}
            \begin{itemize}
                \item \textbf{Windows:} MiKTeX (ligero) o TeX Live.
                \item \textbf{Linux:} TeX Live (\texttt{sudo apt install texlive-full}).
            \end{itemize}
        \item \textbf{Un Editor (IDE):}
            \begin{itemize}
                \item \textbf{Texmaker:} Multiplataforma, muy estable y gratuito.
                \item \textbf{TeXstudio:} Con autocompletado avanzado.
                \item \textbf{VS Code:} Con la extensión "LaTeX Workshop".
            \end{itemize}
    \end{enumerate}
\end{frame}

\begin{frame}[standout]
    ¡Gracias! \\
    ¿Preguntas?
\end{frame}

\section{Referencias Bibliográficas}

\begin{frame}{Referencias en Formato IEEE}
    % Nota: Para mostrar esto realmente, tendrias que cambiar el estilo en el preambulo 
    % a [style=ieee]. Aqui lo mostramos como ejemplo visual.
    \begin{enumerate}
        \item[\small [1]] \small Manual de LaTeX, ``Símbolos Matemáticos'', [En línea]. Disponible en: \url{https://manualdelatex.com/simbolos}.
        \item[\small [2]] CodeCogs, ``Editor Visual de Ecuaciones'', 2026. [En línea]. Disponible en: \url{https://editor.codecogs.com}.
        \item[\small [3]] Tables Generator, ``Generador de tablas para LaTeX'', [En línea]. Disponible en: \url{https://www.tablesgenerator.com}.
    \end{enumerate}
\end{frame}

\begin{frame}{Referencias en Formato APA 7}
    \begin{itemize}
        \item[\quad] \small CodeCogs. (2026). \textit{Editor Visual de Ecuaciones}. Recuperado de \url{https://editor.codecogs.com}
        \item[\quad] \small Manual de LaTeX. (s.f.). \textit{Símbolos Matemáticos}. Recuperado de \url{https://manualdelatex.com/simbolos}
        \item[\quad] \small Tables Generator. (s.f.). \textit{Generador de tablas para LaTeX}. Recuperado de \url{https://www.tablesgenerator.com}
    \end{itemize}
\end{frame}

\end{document}
